%!TEX root = ../../template.tex

\typeout{NT FILE chapter2}%

\chapter{Related Work}
\label{cha:related-work}

In this chapter, we will discuss the architecture and design of the \emph{Semantic Web} as well as analyzing work already done in the fields of \emph{Secure Computations}. The remaining of the chapter is organized as follows: in Section \ref{ch2:sec:semantic-web} we discuss the \emph{Semantic Web} stack, languages and tools; Section \ref{ch2:sec:secure-computations} describes work on \emph{Secure Computation}, with particular attention regarding \gls{HE} and \gls{SE}; finally, in Section \ref{ch2:sec:summary} we summarize and discuss which work is relevant.

\prependtographicspath{{Figures/}}

\section{Semantic Web: A layered architecture}
\label{ch2:sec:semantic-web}

In Section \ref{ch1:sec:context}, we have contextualized the main philosophy and design principles of the \emph{Semantic Web}:
\begin{enumerate}
    \item Usage of \emph{labeled graphs} as the data model where concepts are represented as triples consisting of a subject, predicate and object (with nodes being objects and relations between such objects modelled as the edges);
    \item All data items are identified by \gls{IRI} \autocite{Durst2005-xq};
    \item Usage of ontologies to model semantics within a domain.
\end{enumerate}

\begin{figure}[H]
    \centering
    \includegraphics[width=0.4\textwidth]{chapter2/semantic-web-stack.png}
    \caption{The \emph{Semantic Web} stack \protect\autocite{Wikipedia_contributors_undated-bj}.}
    \label{fig:semantic-web-stack}
\end{figure}


Additionally, the \emph{Semantic Web} can also be seen as a stack of various layers \autocite{Antoniou2012-nj} (Figure \ref{fig:semantic-web-stack}). Most of these layers are standardized due to the collective work of the \gls{W3C}.

Briefly describing this layered architecture, at the root level, we have the \gls{IRI}s and a character set providing the baseline of communication, with distinguishable identification of resources and an agreed encoding.

At the syntax level, we find the \gls{XML} \autocite{W3c_undated-ep}. \gls{XML} supports the definition of unlimited \emph{tags}, that can be placed in the text of a document. These \emph{tags} facilitate the creation of hierarchy and give \emph{structure} to the information.

From structured data we then define the basic data model using \gls{RDF} \autocite{W3c_undated-vq}. With \gls{XML} like syntax, we can write statements about resources, thus embedding  \emph{semantics} within the information, as \emph{metadata}.

\gls{RDFS} \autocite{W3c_undated-mp} is a vocabulary description language to organize \gls{RDF} resources into hierarchies, based on primitives such as \emph{classes} and \emph{properties}. It can be seen as a rudimentary \emph{ontology language} \autocite{Antoniou2012-nj}.

Ontologies are defined with \gls{OWL2} \autocite{W3c_undated-dr}. Like \gls{RDFS}, it is also a vocabulary description language, but provides more expressiveness (e.g. relations between classes, \emph{cardinality}, \emph{equality}, richer properties with characteristics).

\gls{RDF} triples can be interpreted as logical binary predicates, making reasoning possible over \gls{RDF}-based data. This allows us to infer and validate ontologies, and also provide explanations on operations' results. The rules layer is a reference to this specific part of the architecture stack. There exist various languages, such as the \gls{SWRL} \autocite{W3c2004-rn}, that allow us to express rules.

To query and retrieve any data based on \gls{RDF} we use \gls{SPARQL}.

The Proof and Unifying Logic layers include the notion of interconnectivity of terminologies between \emph{knowledge bases}, with explainable proofs to operations.

Only with secure cryptographic schemes can we reach the trust layer. The security layer is still not standardized.

The top layer, references all the applications and interfaces that use this new paradigm of data representation.

In the following of this section we will describe, in more detail, some languages and tools. We will focus on \gls{RDF} in Section \ref{ch2:sec:semantic-web:rdf}; on \gls{RDFS}, in Section \ref{ch2:sec:semantic-web:rdf-schema}; on \gls{OWL2}, in Section \ref{ch2:sec:semantic-web:owl}; on \gls{SPARQL}, in Section \ref{ch2:sec:semantic-web:sparql}; on common tools and software, in Section \ref{ch2:sec:semantic-web:tools}.

\subsection{RDF: Resource Description Framework}
\label{ch2:sec:semantic-web:rdf}

\subsubsection{Overview}
\gls{RDF} \autocite{W3c_undated-vq} is a flexible \emph{domain-independent} data model. The \gls{RDF} basic primitives are \emph{resources}, \emph{properties} and \emph{literals}. We can construct \emph{statements} about resources by organizing these basic primitives in sets of \emph{triples}. A triple is composed of a \emph{subject}, a \emph{predicate} and an \emph{object}. In a statement (triple) about a subject (always a resource) we define the value (object) of a property (predicate), which can be another resource or a literal.

Literals are atomic values (e.g. strings, numbers). Resources and properties, respectively, model the "things" and relationships between such "things". Both are identified using \gls{IRI}s. Resources without an \gls{IRI} are called \emph{blank nodes} (or \emph{bnodes}).

To make a statement $B$ about a statement $A$ we can use a technique called \emph{reification}. Reification consists in: firstly, defining a statement $B$ where the object $objA$ represents the statement $A$; secondly, creating the properties \emph{subject}, \emph{predicate} and \emph{object} mapped from $objA$ to the original subject, predicate (now a resource), and original object of statement $A$ \autocite{Antoniou2012-nj}.

\subsubsection{Syntax}

\paragraph{Graphical} \gls{RDF} being a data model, can be written using different syntaxes. Graphically we can represent statements as a labeled and directed graph. The \emph{nodes} are labelled with the resources' \gls{IRI}s, a literal value or left blank, whereas \emph{arcs} are labelled with the properties' \gls{IRI}s. The arcs are directed from subject to object of a statement. Figure \ref{fig:rdf-grafical-example} provides an example of a \gls{RDF} graph.


\begin{figure}[H]
    \centering
    \includegraphics[width=0.8\textwidth]{chapter2/rdf-example.png}
    \caption{A simple RDF graphical example.}
    \label{fig:rdf-grafical-example}
\end{figure}

\paragraph{Text-Based}


There are multiple text-based syntaxes for \gls{RDF}. Some of the most commonly used are \gls{TURTLE} \autocite{W3c_undated-kq} (Listing\ref{listing-turtle-example}), and its extension for named graphs \emph{TriG} \autocite{W3c_undated-uz}.
In \gls{TURTLE}, \gls{IRI}s are enclosed by angle brackets. The subject, property, and object are written in this order. Statements are finished by a period. Literals are surrounded in quotes with the type of data specified with an \gls{IRI}. The data value and type are separated by \^{}\^{} (if no data type is mentioned, the value is inferred as being a string). We can also define \gls{IRI} abbreviations (using @prefix, followed by the prefix to substitute the \gls{IRI}, a colon, the \gls{IRI} and a period) and simplify statements, for easier readability (e.g.using a semicolon to make more than one statement about the same subject, defining statements with nested blank nodes using square brackets). Comments start with \#.

\lstinputlisting[language=XML, caption={Simple RDF example, written in \gls{TURTLE} syntax.}, label={listing-turtle-example}]{Code/chapter2/turtle-example.xml}

\emph{XML/RDF} \autocite{W3c_undated-hv} is a standardized \gls{XML}-based syntax and \emph{RDFa} \autocite{W3c_undated-qb} permits \gls{RDF} to be embedded in the \gls{HTML}\footnote{We recommend the curious reader to read the W3C recommendations to better learn the syntaxes.}.

\subsection{RDFS: RDF Schema}
\label{ch2:sec:semantic-web:rdf-schema}
%Describe structure, give example extending on the RDF example, enumerate terms and possible syntaxes

\subsubsection{Overview}

\gls{RDF} is a flexible data-model. However, for interconnectivity of different \gls{RDF} knowledge bases, semantics for restricting meaning in a specific domain are needed.

\gls{RDFS} \autocite{W3c_undated-mp} does that by providing vocabulary to describe groups of related resources and the relationships between these resources. This vocabulary is a set of \gls{RDF} resources identified under \url{http://www.w3.org/2000/01/rdf-schema#}{}.

\begin{figure}[H]
    \centering
    \includegraphics[width=0.8\textwidth]{chapter2/rdfs-example.png}
    \caption{A simple RDFS graphical example.}
    \label{fig:rdfs-grafical-example}
\end{figure}

\subsubsection{RDFS primitives}

\emph{Classes} group resources together, the \emph{instances} of the class. \emph{Properties} are defined by specifying to which classes of resource they may apply, using the mechanism of \emph{domain} and \emph{range}. This allows for easy extension, by not having to redefine classes wherever a new property, that applies to an existing class, is needed.

To state that a resource is an instance of a class we use the \emph{rdf:type} property.

The \emph{rdfs:subClassOf} is used to define that all the instances of a class are instances of its \emph{superclasses}: if $A$ rdfs:subClassOf $B$, then instances of $A$ are also instances of $B$. The \emph{rdfs:subPropertyOf} states that instances related by a property are also related by its \emph{superproperties}: if $A$ rdfs:subPropertyOf $B$, then instances related to $A$ are also related to $B$. The properties rdfs:subClassOf and rdfs:subPropertyOf are \emph{transitive}, and there are no restrictions on the number of superclasses or superproperties that, respectively, a class or a property can have.

The \emph{rdfs:domain} property defines that any resource with a given property is an instance of the domain classes that the property applies to. The \emph{rdfs:range} property specifies that the values of the property are instances its range classes.

In tables \ref{tbl:rdfs-classes} and \ref{tbl:rdfs-properties}, we summarize the core \gls{RDFS} classes and properties. In Figure \ref{fig:rdfs-grafical-example}, we extend the \gls{RDF} graphical Example \ref{fig:rdf-grafical-example} with \gls{RDFS}.

\begin{table}[h]
    \begin{center}
        \begin{tabular}{ |l|l| }
            \hline
            \multicolumn{2}{|c|}{\textbf{\gls{RDFS} Core Classes}} \\
            \hline
            \textbf{Class} & \textbf{Description} \\
            \hline
            rdfs:Resource & The class resource, everything. \\
            rdfs:Class & The class of classes. \\
            rdfs:Literal & The class of literal values, e.g. textual strings and integers. \\
            rdf:Property & The class of RDF properties. \\
            rdfs:Statement & The class of RDF statements. \\
            \hline
        \end{tabular}
    \end{center}
    \caption{\gls{RDFS} Core Classes and respective description.}
    \label{tbl:rdfs-classes}
\end{table}

\begin{table}[h]
    \begin{center}
        \resizebox{\textwidth}{!}{\begin{tabular}{ |l|l|l|l| }
            \hline
            \multicolumn{4}{|c|}{\textbf{Core \gls{RDFS} Properties}} \\
            \hline
            \textbf{Property} & \textbf{Description} & \textbf{Domain} & \textbf{Range}\\
            \hline
            \multicolumn{4}{|c|}{\textbf{Relationship Definition}} \\
            \hline
            rdf:type & The subject is an instance of a class. & rdfs:Resource & rdfs:Class \\
            rdfs:subClassOf & The subject is a subclass of a class. & rdfs:Class & rdfs:Class \\
            rdfs:subPropertyOf & The subject is a subproperty of a property. & rdf:Property & rdf:Property \\
            \hline
            \multicolumn{4}{|c|}{\textbf{Property Restriction}} \\
            \hline
            rdfs:domain & A domain of the subject property. & rdf:Property & rdfs:Class \\
            rdfs:range & A range of the subject property. & rdf:Property & rdfs:Class \\
            \hline
            \multicolumn{4}{|c|}{\textbf{Reification}} \\
            \hline
            rdf:subject & The subject of the subject RDF statement. & rdf:Statement & rdfs:Resource \\
            rdf:predicate & The predicate of RDF statements. & rdf:Statement & rdfs:Resource \\
            rdf:object & The object of RDF statements. & rdf:Statement & rdfs:Resource \\
            \hline
            \multicolumn{4}{|c|}{\textbf{Utility}} \\
            \hline
            rdfs:seeAlso & Further information about the subject resource. & rdfs:Resource & rdfs:Resource \\
            rdfs:isDefinedBy & The definition of the subject resource. & rdf:Resource	 & rdfs:Resource \\
            rdfs:comment & A description of the subject resource. & rdf:Resource & rdfs:Literal \\
            rdfs:label & 	A human-readable name for the subject. & rdf:Resource & rdfs:Literal \\
            \hline
        \end{tabular}}
    \end{center}
    \caption{\gls{RDFS} Core Properties and respective description.}
    \label{tbl:rdfs-properties}
\end{table}

\subsection{OWL2: Web Ontology Language}
\label{ch2:sec:semantic-web:owl}

\subsubsection{Overview}

An \emph{ontology} was defined by Studer et al. as "a formal, explicit specification of a shared conceptualisation" \autocite{Studer1998-kp}. An ontology should specify consensual and machine-readable knowledge by means of a formal definition of concepts, where concepts have explicit types and restrictions. Therefore, an \emph{ontology language} should provide a well-defined syntax, formal semantics, convenient and sufficient expressiveness and efficient reasoning support. 

\gls{RDFS} can be considered a very limited \emph{ontology language}. \gls{OWL2} \autocite{W3c_undated-dr} provides a number of additional features (e.g. cardinality constraints, class equivalence, intersection and disjunction). \gls{OWL2} is divided in two sublanguages: \emph{OWL2 Full} \autocite{W3c_undated-ae}, a direct extension of \gls{RDFS} with all \gls{OWL2} semantics, which is undecidable \autocite{Antoniou2012-nj}; \emph{OWL2 DL} \autocite{W3c_undated-xs}, a language with some restrictions on the way \gls{OWL2}, \gls{RDF} and \gls{RDFS} primitives may be used.

The correspondence theorem of \autocite{W3c_undated-xs} states that: given an \emph{OWL2 DL} ontology, inferences drawn using its \emph{Direct Semantics} will still be valid if the ontology is mapped into an RDF graph and interpreted using the \emph{RDF-Based Semantics}, the semantics of \emph{OWL2 Full}.

\gls{OWL2} can be expressed with any valid \gls{RDF} syntax, however there are various specific syntaxes: the Functional-Style Syntax \autocite{W3c_undated-xm}, which is very compact and was the language used in the formal specification of the language; the Manchester Syntax \autocite{W3c_undated-dg}, designed to be human-readable (it is very popular in user interfaces of semantic editors); and the OWL/XML Syntax \autocite{W3c_undated-gr}, which allows easier processing with \gls{XML} tools.

Figure \ref{fig:owl2-structure} details the structure of \gls{OWL2}. Finally, \emph{OWL2} can have different specifications, called profiles \autocite{W3c_undated-jk}, which provide reasoning in polynomial time. The coverage of \gls{OWL2} profiles is not in the scope of this work.

\begin{figure}[H]
    \centering
    \includegraphics[width=\textwidth]{chapter2/OWL2-structure.png}
    \caption{The \gls{OWL2} Structure \protect\autocite{W3c_undated-dr}.}
    \label{fig:owl2-structure}
\end{figure}

\subsubsection{OWL Language}

In \gls{OWL2} the term individual is equivalent to instance in \gls{RDFS}. An \emph{assertion} is a statement about a resource. \emph{Expressions} are a combination of classes, properties and instances. An \emph{axiom} relates such expressions to classes\footnote{For an overview of the language we recommend the reader to consult the reference guide \autocite{W3c_undated-ls}.}. Listing \ref{listing-owl2-example} provides a simple \gls{OWL2} example.

\paragraph{Individual Assertions} In \gls{OWL2} \emph{assertions} about individual resources can be used to define their type, class membership, properties and even provide ways to define identity and negative statements. More specifically on the identity statements, because \gls{OWL2} has an open world assumption, we can never differentiate two individuals just because they have different \gls{IRI}s. To do so we use owl:differentFrom. To explicitly state negative facts we use owl:NegativePropertyAssertion \autocite{Antoniou2012-nj}.

\lstinputlisting[language=XML, caption={Simple \gls{OWL2} example, written in \gls{TURTLE} syntax.}, label={listing-owl2-example}]{Code/chapter2/owl2-example.xml}

\paragraph{Classes} In \gls{OWL2} every resource that is of type owl:Class is a class. The two classes defined by default are \emph{owl:Thing} and \emph{owl:Nothing}. Respectively, these two specify the superclass of all classes and the subclass of all classes: every individual is a member of owl:Thing and no individual can be member of owl:Nothing. This is very useful for reasoning, because inconsistency is detected when classes are equivalent to owl:Nothing (have no individuals in their membership).

\paragraph{Class Axioms} In \gls{OWL2} we can relate \emph{expressions} to classes using \emph{axioms}. These can define:

\begin{itemize}
    \item\emph{Subclass Relations}, with equivalent usage like in \gls{RDFS};
    \item\emph{Class Equivalence}, in equivalent classes every individual must be a member of each class;
    \item\emph{Enumerations}, a type of class definition via explicitly stating all individuals who belong to the class membership;
    \item\emph{Intersection}, in a class described as the intersection of various classes, all of its members must be also individuals belonging to the membership of each of those classes;
    \item \emph{Union and Disjoint Union}, in a class described as the union of various classes, every individual belonging to its membership must also belong to the membership of at least one of those classes. In the particular case of disjoint union, the intersection of the classes is equivalent to owl:Nothing;
    \item\emph{Disjoint Classes}, in disjoint classes no individual can be a member of both classes (the intersection of the classes is equivalent to owl:Nothing);
    \item\emph{Complement}, complement classes contain all the individuals such that the union of both classes is equivalent to owl:Thing.
\end{itemize}


\paragraph{Properties} \gls{OWL2} has three types of properties: \emph{object} properties, which relate individuals to other individuals; \emph{datatype} properties, which relate individuals to literals; and \emph{annotation} properties; which are used to specify human-readable labels, comments and explanations. 


The \emph{top} and \emph{bottom} properties specify, respectively, properties that relate to all and no individuals: owl:topObjectProperty is the superproperty of all owl:ObjectProperties and owl:bottomObjectProperty relates to no individual; owl:topDataProperty is the superproperty of all owl:DatatypeProperties, relating an individual to any possible literal value, and owl:bottomDataProperty relates no individual to any literal value.

Properties can have different characteristics:
\begin{itemize}
    \item \emph{Transitive Properties}: if $a$ relates to $b$ via property $p$ and $c$ relates to $a$ via property $p$, then $c$ relates to  $b$ via property $p$;
    \item \emph{Symmetric Properties}: if $a$ relates to $b$ via property $p$, then inverse is true;
    \item \emph{Asymmetric Properties}: if $a$ relates to $b$ via property $p$, then the inverse is $\bot$;
    \item \emph{Functional Properties}: if for each instance the property can only have one unique value;
    \item \emph{Inverse-functional Properties}: if the object of a property statement relates to a single subject;
    \item \emph{Reflexive Properties}: if every individual relates to itself via the property;
    \item \emph{Irreflexive Properties}: if no individual relates to itself via the property.
\end{itemize}

\paragraph{Property Axioms} Just like class axioms, we can specify \emph{property axioms} to describe the way properties relate to classes and other properties. These allow us to define:

\begin{itemize}
    \item \emph{Domain and Range}: used to determine class membership of individuals;
    \item \emph{Inverse Relationship}: used to specify the inverse of the property;
    \item \emph{Equivalence}: used to specify equivalent properties. Individuals related by one property will always relate by all equivalent properties;
    \item \emph{Disjointness}: used to specify disjoint properties, where individuals related by one property can never be related by the other;
    \item \emph{Property Chains}: used to aggregate connected properties under another property.
\end{itemize}
 

\paragraph{More Granular Class Axioms} We can write more granular class definitions in \gls{OWL2} than by simply defining class axioms. This is achieved by combining class and property axioms. It is done by defining a class axiom that relates to an anonymous class (owl:Restriction) which contains axioms on its properties. The anonymous class definition captures all individuals that satisfy such restrictions. As such we can define:
\begin{itemize}
    \item \emph{Universal and Existential Restrictions}: which the members of the class must satisfy. The class is defined as a subclass of an anonymous class with property restrictions on, respectively, owl:allValuesFrom and owl:someValuesFrom;      
    \item \emph{Necessary and Sufficient Conditions}: relating a class to a property, using rdfs:subClass\-Of, does not provide sufficient information to conclude whether an individual belongs to the class membership or not (it only states necessary, not sufficient, conditions for class membership). With owl:equivalentClass, it becomes clear that an individual belongs to the class membership, if it satisfies the restriction (the conditions are necessary and sufficient).
    \item \emph{Value, Cardinality, Data Range and Datatype Restrictions} on the anonymous class: limiting the value, cardinality of the value, range of possible values and the datatype of the property value relating to members of the superclass; 
    \item \emph{Self Restrictions}: on the individuals captured by the anonymous class.

\end{itemize}

\subsubsection{OWL2 DL}

Contrary to \emph{OWL2 Full}, \emph{OWL2 DL}, being less expressive, retains \emph{decidability}. This is achieved by imposition of various restrictions \autocite{Antoniou2012-nj}: 

\begin{itemize}
    \item The application of \gls{OWL2} primitives to each other is not allowed; only classes and non-literals resources may be defined, with all classes being individuals of owl:Class;
    \item Properties are disjoint individuals of either owl:ObjectProperty or owl:DatatypeProperty;
    \item Annotation properties are ignored by reasoners;
    \item Properties, for which range includes non-literal resources, are separated from those which relate to literal values; 
    \item A resource can only be a class, property, or instance at the same time. Classes, properties or instances with the equal names will be treated as distinct (a mechanism called punning);
    \item Only functional properties can be assigned to datatype properties;
    \item Only object properties can have an inverse;
    \item Property chains can only involve object properties;
    \item Self restrictions on datatype properties are not allowed.
\end{itemize}

\subsection{SPARQL}
\label{ch2:sec:semantic-web:sparql}

\subsubsection{Overview}

In the previous section we have discussed the structure of \gls{OWL2}. \gls{OWL2} documents have a direct mapping to \gls{RDF} documents. As such, we can publish and query \gls{OWL2} ontologies using the standard approach for \gls{RDF} data. 

\gls{RDF} triples are stored in \emph{triple stores}, specialized databases for this type of data. The underlying architecture of such databases varies, from relational to NoSQL approaches \autocite{Cure2014-vm, Modoni2014-cz}. The recommended query language to use is \gls{SPARQL} \autocite{W3c2013-hd}. It provides a range of operations over published \gls{RDF} graphs, accessible via SPARQL endpoints: from read queries \autocite{W3c2008-vf} to write operations \autocite{W3c_undated-tw, W3c2013-rr}. 

\subsubsection{SPARQL Queries}
In this section we will mathematically define core components of \emph{SPARQL}, briefly discuss the basis of functioning for this type of queries, explain how they are evaluated, given the semantics of \emph{graph patterns}, and cover their syntax. 

\begin{definition}[\textbf{RDF dataset}] Let $I$, $B$, $L$ be disjoint sets of \gls{IRI}s, blank nodes, and literals, respectively. \gls{RDF} data is a set of subject, predicate and object triples $t = (s,p,o)$, where $s \in I \cup B$, $p \in I$ and $o \in I \cup B \cup L$. 

Let $V$ be a set of nodes formed by subjects and objects and $E$ be a set of edges, labeled by the corresponding predicate. A set of \gls{RDF} triples can be represented as a directed, edge-labeled, graph $G = (V, E)$. An \gls{RDF} dataset is a set of \emph{named} graphs and the \emph{default} unnamed graph. We will define such dataset as $D = \{G, \gamma\}$, where $\gamma$ is a function which assigns a named graph $\gamma(i)$ to each $i$ in a finite set $dom(\gamma) \in I \cup B$ \autocite{Picalausa2011-ca, noauthor_undated-ev}.
\end{definition}

\begin{definition}[\textbf{SPARQL Query}] Let $T = \{\textrm{SELECT},\; \textrm{ASK},\; \textrm{CONSTRUCT},\;\textrm{DESCRIBE}\}$, a \gls{SPARQL} query can be defined as a tuple $q = (t_q,d,m,P_q)$ where $t_q \in T$ is the \emph{query type}, $d \in I \cup B$ is an optional term identifying the \gls{RDF} graph to be queried, $m$ is a set of \emph{solution modifiers} and $P_q$ is a \emph{graph pattern}.               
\end{definition}

\begin{definition}[\textbf{Graph pattern}]\label{ch2:definition:sparql-graph-pattern}
Let $V = \{?v_0,...,?v_ n\}$, where $n \in \mathbb{Z}^{+}$, be a set of variables, disjoint from $I$, $B$, and $L$. A triple pattern is a tuple $tp = (s', p', o')$ such that $s' \in I \cup V$, $p' \in I \cup V$ and $o' \in I \cup L \cup V$. A set of triple patterns is usually referred as a \gls{BGP}. A \emph{property path pattern} $pp$ can be inductively defined as an element of $(I \cup B \cup V) \times pp \times (I \cup B \cup L \cup V)$. Finally, a \emph{graph pattern} is an expression specified by the following grammar:

\begin{equation}
    \begin{split} 
    P ::= \;\; & tp \;\;\vert\;\; pp \;\;\vert\;\; q \;\;\vert\;\; P_1\; op \; P_2 \;\;\vert\;\; P\; \textrm{FILTER}\; R \;\;\vert\;\; \textrm{Graph}\; i\; P \;\;\vert\;\; \textrm{Graph}\; v?\; P
    \end{split}
\end{equation}


Here $tp$ is a triple pattern, $pp$ is a \emph{property path pattern}, $q$ is a subquery, $op = \{ \textrm{UNION}, \allowbreak \textrm{OPTIONAL}, \; \textrm{MINUS}, \; . \}$ represents a set operation between two graph patterns, $R$ is a \gls{SPARQL} filter constraint, $i \in I$, $?v \in V$ and the keyword \emph{Graph} is used to specify which dataset to match the pattern against \footnote{For a detailed explanation of the syntax and semantics, we recommend the reader to consult section 17 of \autocite{W3c2008-vf}}. To remove ambiguities parentheses can be added to a pattern. Finally, we will define the set of variables that occur in a pattern $P$ as $vars(P)$ \autocite{Picalausa2011-ca}.
\end{definition}

\paragraph{How Do SPARQL Queries Work?} An isomorphism of two graphs is an adjacency preserving bijection between their vertex sets \autocite{McKay2014-zb}.

The basis of \gls{SPARQL} queries is finding matches with \emph{graph patterns} in the \gls{RDF} dataset \autocite{Ladwig2013-xp, Cure2014-vm}. This problem is equivalent to the \emph{subgraph isomorphism problem}, a generalization of the graph isomorphism problem: given two graphs $G_1$ and $G_2$, we try to find if $G_1$ contains a subgraph that is isomorphic to $G_2$. 

In the context of a \gls{SPARQL} query $q$, $G_1$ is a published \gls{RDF} graph identified by $d$ and $G_2$ corresponds to a \emph{graph pattern} $P_q$. From this match we obtain a set of \emph{bindings}: a multiset of mappings between $vars(P)$ and data within $G_d$. The final result of a query will also depend on $t_q$, the \emph{query type}, and $m$, the specified \emph{solution modifiers} (i.e. aggregation, grouping, sorting, duplicate removal). 

\begin{definition}[\textbf{Semantics of SPARQL graph patterns}] A \emph{mapping} $\mu$ is a partial function $\mu : \; V \to \; I \cup B \cup L$ that assigns values to a finite set of variables. The domain of $\mu$, denoted by $dom(\mu)$, is the subset of $V$ where $\mu$ is defined. 

Two \emph{mappings} $\mu_1$ and $\mu_2$ are compatible, denoted by $\mu_1 \sim \mu_2$, if $\mu_1(?v) = \mu_2(?v)$ for all common variables $?v \in dom(\mu_1) \cap dom(\mu_2)$.

Abusing notation, we define $t = \mu(tp)$ as the triple obtained after replacing the variables in the \emph{triple pattern} $tp$ according to $\mu$ \autocite{Letelier2013-rs, Picalausa2011-ca}. Finally, let $\mu \vDash R$ represent a mapping that satisfies a filter condition $R$.

We define the \emph{join}, \emph{union}, \emph{difference}, and \emph{left outer join} operations on two sets of mappings $\Omega_1$ and $\Omega_2$ as follows:

\begin{equation}
    \begin{split} 
    & \Omega_1 \Join \; \Omega_2 = \{ \mu_1 \cup \mu_2 \;\vert\; \mu_1 \in \Omega_1 ,\; \mu_2 \in \Omega_2,\; \mu_1 \sim \mu_2 \} \\ 
    & \Omega_1 \cup \; \Omega_2 = \{ \mu \;\vert\; \mu \in \Omega_1 \vee \mu \in \Omega_2 \} \\
    & \Omega_1 - \; \Omega_2 = \{ \mu_1 \in \Omega_1 \;\vert\; \forall \mu_2 \in \Omega_2,\; \mu_1 \nsim \mu_2 \} \\
    & \Omega_1 \leftouterjoin \; \Omega_2 = (\Omega_1 \Join \Omega_2) \; \cup \; (\Omega_1 - \; \Omega_2) \\
    \end{split}
\end{equation}

The semantic evaluation of $P$ over a dataset $D$ is a set of mappings $\llbracket P \rrbracket_D$. Inductively, we define it as:

\begin{equation}
    \begin{split} 
    & \llbracket tp \rrbracket_D = \{ \mu \;\vert\; dom(\mu) = vars(tp) \wedge \mu(tp) \in G_D \}\\
    & \llbracket pp \rrbracket_D = \{ \mu \;\vert\; dom(\mu) = vars(pp) \wedge \mu(pp) \in G_D \}\\
    & \llbracket q \rrbracket_D = \{ \mu \;\vert\; dom(\mu) = vars(P_q) \wedge \mu(P_q) \in G_D \}\\
    & \llbracket P_1 \;.\; P_2 \rrbracket_D = \llbracket P_1 \rrbracket_D \Join \llbracket P_2 \rrbracket_D \\
    & \llbracket P_1 \;\textrm{UNION}\; P_2 \rrbracket_D = \llbracket P_1 \rrbracket_D \cup \;\llbracket P_2 \rrbracket_D \\
    & \llbracket P_1 \;\textrm{OPTIONAL}\; P_2 \rrbracket_D = \llbracket P_1 \rrbracket_D \leftouterjoin \; \llbracket P_2 \rrbracket_D \\
    & \llbracket P_1 \;\textrm{MINUS}\; P_2 \rrbracket_D = \llbracket P_1 \rrbracket_D - \; \llbracket P_2 \rrbracket_D \\
    & \llbracket P_1 \;\textrm{FILTER}\; R \rrbracket_D = \{ \mu \in \llbracket P \rrbracket_D \;\vert\; \mu \vDash R \} \\
    & \llbracket \textrm{Graph} \;i\; P\rrbracket_D = \{ \mu \in \llbracket P \rrbracket_{(\gamma(i), \gamma)} \;\vert\; i \in dom(\gamma) \}\\
    & \llbracket \textrm{Graph} \;?v\; P\rrbracket_D = \bigcup_{i \; \in \; dom(\gamma)} \{ ?v \to i \} \Join \llbracket P \rrbracket_{(\gamma(i), \gamma)} \\
    \end{split}
\end{equation}
\end{definition}

\paragraph{Syntax} \gls{SPARQL} queries have a similar structure to the \gls{SQL} queries: we define the operation to execute using a keyword, followed by the projection of variables to be shown in the result, the keyword \emph{WHERE}, the \emph{graph pattern}, enclosed in brackets, and, optionally, some modifiers of the result (i.e. \emph{ORDER BY} to sort results and \emph{DISTINCT} to remove duplicates). Additionally, we may also specify the \gls{RDF} graph to consider for the query evaluation with the keywords \emph{FROM} or \emph{FROM NAMED}. Furthermore, the \gls{SPARQL} syntax is very similar to \gls{TURTLE}, so we can specify prefixes for namespaces. 
The basic \emph{\gls{SPARQL} 1.1 Query Language} \autocite{W3c2008-vf} provides the following query types: 

\begin{itemize}
    \item \emph{SELECT}: used to extract values from a \gls{RDF} graph. The result is a set of bindings;
    \item \emph{ASK}: returns a boolean, true if there was a match between the \emph{graph pattern} and the \gls{RDF} graph, false otherwise; 
    \item \emph{CONSTRUCT}: we specify two \emph{graph patterns}, one to fetch \emph{bindings} and another to use as a template for constructing new triples, from the values in the \emph{bindings}. The final result is the set of newly generated triples.
    \item \emph{DESCRIBE}: results depend on implementation \autocite{Picalausa2011-ca}, but usually returns a set of RDF triples describing the matching patterns. 
\end{itemize}

In the following Listings, we present a few query examples. Listing \ref{listing-sparql-example} shows a query with a conjunction, Listing \ref{listing-sparql-modifiers-example} specifies a query with filters and result modifiers and Listing \ref{listing-sparql-complex-example} shows a query with result modifiers and a disjunction.

\pagebreak

\lstinputlisting[language=XML, caption={\gls{SPARQL} query with a conjuntion pattern}, label={listing-sparql-example}]{Code/chapter2/sparql-example.xml}

\lstinputlisting[language=XML, caption={\gls{SPARQL} query with filters and modifiers}, label={listing-sparql-modifiers-example}]{Code/chapter2/sparql-modifiers-example.xml}

\lstinputlisting[language=XML, caption={\gls{SPARQL} query with a disjunction pattern}, label={listing-sparql-complex-example}]{Code/chapter2/sparql-complex-example.xml}

Besides read operations, \emph{\gls{SPARQL} Update} \autocite{W3c_undated-tw}, provides additional options for inserting, publishing and deleting data in triplestores:

\begin{itemize}
    \item \emph{INSERT}/\emph{DELETE}: these operations add and delete, respectively, newly generated triples, based on a template, into a destination graph (if it does not exist it should be created). Like in a \emph{CONSTRUCT} query we specify two \emph{graph patterns}, one to fetch bindings and another to use as a template for constructing triples. We can join both operations, in this situation we specify three \emph{graph patterns}, two to use as templates to generate the triples to be inserted or deleted and a third one to fetch bindings.
    \item \emph{INSERT}/\emph{DELETE DATA}: similar to their pattern-based variants, but in these operations we specify concrete data inline.
    \item \emph{LOAD}: used to read the contents of a document that represents a graph and saves it into a destination graph in a triplestore.
    \item \emph{CLEAR}: this operation removes all the triples in the specified graphs. 
\end{itemize}

\subsubsection{Preserving Reasoning}

In a \emph{knowledge base} we can distinguish two sets of statements: the \emph{TBox} and the \emph{ABox}. The \emph{TBox} set contains terminology data (e.g. ontologies) which we can define using \gls{RDFS} or \gls{OWL2}. The \emph{ABox} set refers to assertions.

\gls{SPARQL} was not designed for vocabularies like \gls{RDFS} or \gls{OWL2}. By default, a \gls{SPARQL} engine solely considers explicit data, as if it were at the \gls{RDF} layer. Query evaluations interpret all data as assertions, completely ignoring the reasoning capabilities of \gls{RDFS} or \gls{OWL2}. This results in answers that might be incomplete, not containing inferred statements (only derivable from the \emph{ABox} set, using the \emph{TBox} data).

Two popular techniques for preserving reasoning capabilities are \emph{materialization} and \emph{query rewriting}.

\paragraph{Materialization} In this approach, we derive any implicit information and store it alongside the original data. This solution insures fast query evaluation that requires no inference, but suffers from high storage costs. Additionally, write operations are more complex, due to the problem of \emph{truth maintenance}. After each update, we need to maintain the materialized data in a consistent state. 

Insertions are fairly simple, we just need to compute new implicit triples from the data already saved in the system.

Deletions, however, are more complex. We need to compute the set of triples to be removed and assure that no triple in it is entailed without the removed assertions. We usually find three types of implementations for this operation: to always recompute the materialization after a deletion; removing the inferred triples and recomputing the closure of the remaining; maintaining a data-structure with the state of deletions. Neither of these options is perfect, the first two are not efficient for databases with high volume of write operations and the last one incurs in increased costs (associated with maintaining an additional structure).

\paragraph{Query Rewriting} This technique consists in deriving a complete answer set at query time. Before evaluating a query, the system modifies it according to the statements in the \emph{TBox}. Query evaluation is slower: rewriting the queries requires communication with a reasoner and the reformulated queries tend to be larger.

Contrary to the first solution, there is no problem with \emph{truth maintenance}. Write operations can be processed without any additional computations. 


\subsection{Industry solutions}
\label{ch2:sec:semantic-web:tools}

\subsubsection{Development Software \& Tools} 
To developed and maintain ontologies, professionals use advanced semantic editors. \emph{Protégé} \autocite{Musen2015-kh} is a free, open-source, and the usual choice. It provides a rich ontology editing environment with full support for \gls{OWL2}, available in a desktop version or on the web \autocite{Tudorache2009-ja}. Many of its functionalities are based on plugins: automatic verification of errors and consistency in the design using reasoners like \emph{Pellet} \autocite{Sirin2007-lo}; debugging tools; query answering; collaborative development \autocite{Tudorache2008-tu}; graphic visualization of data and knowledge. Additionally, for developers, many \gls{API}s \autocite{Owlcs_undated-im} or full framework solutions, such as \emph{Apache Jena} \autocite{The_Apache_Software_Foundation_undated-wk}, are available for building \emph{Semantic Web} applications.

\subsubsection{Triplestores} 

To store \gls{RDF} data, popular solutions are: \emph{MarkLogic} \autocite{MarkLogic2020-pb}, \emph{Virtuoso} \autocite{OpenLink_Software_undated-km}, \emph{Amazon Neptune} \autocite{Amazon_Web_Services_Inc_or_its_affiliates_undated-fj}, \emph{GraphDB} \autocite{Ontotext_undated-mj}, \emph{Apache Jena - TDB} \autocite{The_Apache_Software_Foundation_undated-mg}. While \emph{Apache Jena - TDB} is primarily a \gls{RDF} triplestore, the rest are \emph{multi-model} database systems. 

To test their capability to handle \gls{SPARQL} queries, the \gls{BSBM} is a standard benchmark \autocite{Chris_undated-hq}. The \gls{LUBM} is another benchmark, which is tailored for \gls{OWL2} data \autocite{Guo2005-ja}. It provides synthetic data generation for a single realistic ontology.

\section{Secure Computations}
\label{ch2:sec:secure-computations}
In Section \ref{ch1:sec:context}, we have contextualized the motivations behind the need for secure computation when sharing ontologies over the cloud or other distributed networks.

In this setting it is common to consider the \emph{honest-but-curious}, or \emph{passive}, adversary model. Such an adversary will follow the protocol properly but keep a record of all intermediate computations \autocite{Goldreich2001-sr} in an attempt to learn private information while observing views of the protocol execution \autocite{Evans2018-it}. Security concerns should take into account the operation history, data stored in the server and also the memory access pattern, as the adversary will have access to all of this information.

The most common approach to secure data in outsourced servers is \emph{Encryption at Rest}. In this approach we encrypt all data that will be stored in the untrusted server. To perform computations we need to download the full dataset, decrypt it, and only then can we execute operations on the data. This is too inflexible and computationally expensive. 

Another approach, using \gls{ORAM} \autocite{Goldreich1996-yb,} fits well against this adversary model, as it allows for the concealment of memory access patterns by shuffling and continuously re-encrypting data on memory, but it is still impractical.

Finally, when used exclusively, \gls{HE} \autocite{Acar2018-mp, Fontaine2007-sv} schemes are not efficient enough to be feasible in real-world applications \autocite{Bosch2014-eb}. However, they can be applied in conjunction with other security schemes, such as in some \gls{SE} schemes \autocite{Bosch2014-eb, Chase2010-xx}. For this reason, we will analyze related work in the fields of \gls{HE}, in Section \ref{ch2:sec:homomorphic-encryption}, and \gls{SE}, in Section \ref{ch2:sec:searchable-encryption}. In Section \ref{ch2:sec:property-preserving-encryption}, we will cover property-preserving encryption, as it is useful for the implementation of equality joins and sorting operations. Lastly, in Section \ref{ch2:sec:blinding-techniques}, we will present two secret sharing techniques, used for the obfuscation of values shared between two parties.


\subsection{Homomorphic Encryption}
\label{ch2:sec:homomorphic-encryption}

\subsubsection{Overview}

The term homomorphism, in security, was used for the first time by Rivest et al. in 1978 \autocite{Rivest1978-vt}.
\glsxtrlong{HE} is a cryptographic method based on the concept of algebra homomorphism, where both decryption and encryption are seen as homomorphisms of the relations between the plaintext and the ciphertext \autocite{Acar2018-mp, Fontaine2007-sv}. With \gls{HE}, computations done over encrypted data, are preserved after decryption of the ciphertext. 

\begin{definition}[\textbf{Homomorphic Encryption}]\label{ch2:definition:he}
    An encryption scheme is homomorphic over operation $\bullet_1$ if it satisfies the following equation, where $E$ is the encryption algorithm, $\bullet_1$ and $\bullet_2$ are arbitrary operations and $M$ is the set of all possible messages \autocite{Acar2018-mp, Franz2011-sm}: 

\begin{equation}\label{ch2:sec:homomorphic-encryption:def-eq}
    E(m_{1}) \, \bullet_2 \, E(m_{2}) = E(m_{1} \bullet_1 m_{2}), \forall m_{1}, m_{2} \in M
\end{equation}

\end{definition}

\subsubsection{Types of HE schemes}

Some \gls{HE} schemes preserve additive properties \autocite{Paillier1999-js, Damgard2007-km, Damgard2009-my}, others maintain multiplicative properties \autocite{Rivest1978-mu, Elgamal1985-rk}, and some conserve both additive and multiplicative properties \autocite{Gentry2009-zv}. The \gls{HE} schemes that preserve both properties allow for any arbitrary operation, because boolean circuits can be represented using only \emph{XOR} (addition) and \emph{AND} (multiplication) gates \autocite{Acar2018-mp}. 

\gls{HE} is classified in four main types: 
\begin{itemize}
    \item \emph{\glsxtrlong{FHE}}: which supports any number of operations with no restriction on the number of operations; 
    \item \emph{\glsxtrlong{SHE}}: which permits some types of operations, a limited number of times;
    \item \emph{\glsxtrlong{PHE}}: that allows for one type of operation with no restriction on the number of times;
    \item \emph{\glsxtrlong{LFHE}} \autocite{Brakerski2014-nm}: which supports any type of operations for a limited, and expected, number of $L$ times (referred to as depth-$L$).
\end{itemize}

\begin{definition}[\textbf{Additive Homomorphic Encryption}]\label{ch2:definition:he-additive} An encryption scheme is as an additive homomorphic scheme if Equation \ref{ch2:sec:homomorphic-encryption:def-eq} holds for the $+$ operator. In this thesis, we will uniquely consider schemes that respect this condition, when $\bullet_1$ is an addition and $\bullet_2$ is a multiplication:

\begin{equation}\label{ch2:sec:homomorphic-encryption:def-eq-additive}
    E(m_{1}) \, * \, E(m_{2}) = E(m_{1} + m_{2})
\end{equation}

\paragraph{Multiplication with Constants} Multiplication by a constant $k > 0$ can be performed on an encrypted value, by raising it to the power of $k$. This is equivalent to performing $k - 1$ multiplications:

\begin{equation}\label{ch2:sec:homomorphic-encryption:def-eq-additive-mul-positive}
    E(m)^k = \underbrace{E(m) \; * \; ... \; *  \; E(m)}_{k \text{ times}} = E(\underbrace{m \; + \; ... \; + \; m}_{k \text{ times}}) = E(k * m)
\end{equation}

Considering that the plaintext space is $M \in \mathbb{Z}_n$ and that messages are small in comparison to the modulus $n$ (so that reductions modulo $n$ never occur), we can represent a negative constant $c$ as $n - c$. Multiplication by a negative value is achieved by raising the encrypted value to the power of $n - c$ (an overflow of module $n$ will occur)\autocite{Franz2011-sm}:

\begin{equation}\label{ch2:sec:homomorphic-encryption:def-eq-additive-mul-negative}
    \begin{split}
        E(m)^{(n - c)} &=  E((n - c)* m \; \textrm{mod} \; n) \\ 
        &=  E(n * m - c * m \; \textrm{mod} \; n)\\ 
        &=  E(-c * m)
    \end{split}
\end{equation}

\paragraph{Subtraction} Given two messages $m_1$ and $m_2$, subtraction is computed by:

\begin{equation}\label{ch2:sec:homomorphic-encryption:def-eq-additive-subtraciton}
    \begin{split}
        E(m_1 - m_2) &=  E(m_1 + (-1) * m_2) \\ 
        &=  E(m_1) * E((-1) * m_2) \\ 
        &=  E(m_1) * E(m_2)^{-1}
    \end{split}
\end{equation}

\end{definition}

\subsection{Property-Preserving Encryption}
\label{ch2:sec:property-preserving-encryption}

\subsubsection{Overview}

\glsxtrlong{HE} allows arbitrary computations to be executed on encrypted data while not leaking any information about the original plaintext. For operations such as equality joins, where we require equality checks, we do need to preserve (leak) certain properties from the original plaintexts. This type of encryption is known as property-preserving encryption.

\subsubsection{Deterministic Encryption} 

Deterministic encryption is type of property-preserving encryption scheme. Contrary to a probabilistic scheme, it will always generate the same ciphertext for each plaintext. This allows equality checks to be performed on the encrypted data. Naturally, its security guarantees are weakened. Deterministic ciphertexts are vulnerable to statistical analysis attacks.

\subsubsection{Order-Preserving Encryption} 

\gls{OPE} schemes preserve the order guarantees of the plaintext data. Naturally, the order relations between the ciphertext are leaked. Some implementations can leak as much as half of the data bits of the ciphertexts \autocite{Popa2012-bw}.

\subsection{Blinding Techniques}
\label{ch2:sec:blinding-techniques}

\subsubsection{Overview} In this thesis we will consider two types of secret sharing techniques, \emph{additive} and \emph{multiplicative blinding}. These are used to obfuscate values shared between parties (e.g. during the execution of a protocol). 

In the following sections we will present how these techniques are implemented. As in \autocite{Franz2011-sm}, we will consider that all computation are done in a modular group $\mathbb{Z}_n$. Additionally, we will denote $P(X = x)$ as the probability that a random variable $X \in \mathbb{Z}_n$ takes on the value $x$.

\subsubsection{Additive Blinding} 

Let $x$ be a fixed value in $\mathbb{Z}_n$. To obfuscate $x$, firstly, we randomly sample $r \in \mathbb{Z}_n$, where the distribution of $r$ is the discrete uniform distribution, thus $P(R = r) = \frac{1}{n}, \forall r \in \mathbb{Z}_n$. After sampling $r$, we compute $y = r + x \; \textrm{mod} \; n$. 

This method generates a uniformly random element $y \in \mathbb{Z}_n$ which does not reveal any information about the value $x$ nor the distribution of $X$.

\subsubsection{Multiplicative Blinding}

We will now assume that $X$ takes on values from a field $\mathbb{Z}_p$, where $p$ is a prime number. Multiplicative blinding is usually used when we want to obfuscate a value $x$, where $x \neq 0$ (if $x$ where to be 0, we would only be able to produce 0 via multiplication).

Let $x \in \mathbb{Z}_p$, and $r \in \mathbb{Z}_p \backslash \{0\}$, we obfuscate $x$ by computing $y = r + x \; \textrm{mod} \; p$. The value of $y$ is equal to 0 if and only if $x = 0$. If $x \neq 0$, then $y$ is a uniformly random element of $\mathbb{Z}_p \backslash \{0\}$ and does not reveal any information about $x$ (a proof is presented in \autocite{Franz2011-sm}). 

\subsection{Searchable Encryption}
\label{ch2:sec:searchable-encryption}

\subsubsection{Overview}
\glsxtrlong{SE} uses classical encryption schemes to perform encrypted search queries. While the first approach implemented it via a searchable ciphertext \autocite{Song2000-hu}, most modern schemes do it using encrypted searchable indexes. \glsxtrlong{SE} can be divided in \gls{SSE}, when it is based on symmetric key cryptography, and \gls{PEKS}, when public key algorithms are used. We will focus on index-based \gls{SE} schemes, more specifically on index-based \gls{SSE} ones, due to their higher efficiency compared to \gls{PEKS} approaches.


%Generic model

\subsubsection{General Index-Based Model}

A general model, for the index-based \gls{SE} schemes was defined in \autocite{Bosch2014-eb}. Given a database $DB=[m_{0},...,m_{n}]$ consisting of a set of $n$ plaintext documents, we build an encrypted searchable index $I=[i_0,...,i_n]$, for the database $DB$, using a \emph{BuildIndex} function. This function extracts some data items (e.g, keywords in the documents) $W=[w_{0},...,w_{m}]$ from the documents and encrypts them under key $K$. The documents are also encrypted, generating a collection of ciphertexts $C=[c_1,...,c_n]$, usually with a different key $K'$. The client then stores the following information in an \emph{honest-but-curious} untrusted server:

\begin{equation}
    \begin{split}
    &I = BuildIndex_{K}(DB, \; W) \\
    &C = E_{K'}(DB) 
    \end{split}
\end{equation}

To perform searches, the client generates a trapdoor $T = Trapdoor_{K}(p)$, for a predicate $p$ (e.g. a keyword $w$). With trapdoor $T$, the server performs the search over $I$, using a $Search$ algorithm, by checking for an encrypted index $i = Search(I,T)$. All documents referenced by index $i$ (which satisfy predicate $p$) are returned to the client. In the client, these are decrypted using key $K$. If no match is found, the server returns $\bot$.

\begin{equation}
    \begin{split}
    &T = Trapdoor_{K}(p) \\
    &\{ C_i = [c_{i,0}, ..., c_{i,n}], \; \bot\} = Search(I, \; T) \\
    &M_p = [m_{p,0},...,m_{p,n}] = D_{K'}(C_i) 
    \end{split}
\end{equation}

\subsubsection{Client/Server Model}

In \gls{SE} schemes, one or various clients, the \emph{writers} store encrypted data on a server. To issue a query request to the server, one or multiple clients, the  \emph{readers}, generate trapdoors. There are four possible \gls{SE} scheme architectures: \gls{S/S}; \gls{S/M}; \gls{M/S}; \gls{M/M}.

The \gls{S/S} architecture is usually used for \emph{data outsourcing} and \gls{M/S}, \gls{S/M} and \gls{M/M} employed in \emph{data sharing} scenarios \autocite{Bosch2014-eb}. Symmetric key schemes are typically \gls{S/S}, while public key schemes \gls{M/S} (multiple users can write using a \emph{public key}, but only the owner of the \emph{private} key can generate trapdoors).
Extension of the single reader schemes to multiple reader schemes can be done with \emph{key sharing} or with the usage of \emph{proxy re-encryption}. 

In proxy re-encryption cryptography, a proxy is used to re-encrypt a ciphertext, encrypted using key $K$, and generates a new ciphertext, that is decrypted using key $K'$, where $K \neq K'$ and the decrypted plaintext remains unchanged. Proxy re-encryption is usually employed with public key cryptography \autocite{Blaze1998-pz, Ivan2003-uu, Dong2008-qw}, but there are some schemes that support the symmetric key setting \autocite{Syalim2011-fb, Sakurai2017-ui, Peng2017-ia}.

\subsubsection{Information Leakage}

Almost all \gls{SE} techniques have privacy issues because they leak some information \autocite{Bosch2014-eb}. This leakage can be \emph{index} information, the \emph{search pattern}, or \emph{access pattern} during execution:

\begin{itemize}
    \item \emph{Index Information}: refers to the information about the index keywords. It may include the number of keywords per document or database, the total number of documents, document similarity, their length and metadata;
    \item \emph{Search Pattern}: includes the information about keywords, and predicates, that is inferred by correlating equal results from different searches. The search pattern can be targeted using statistical analysis;
    \item \emph{Access Pattern}: includes the information that may be derived from the query results, based on the memory accesses needed to retrieve the documents. (e.g. one can infer whether a query is more restrictive than another, based on the total number of documents retrieved; or if some files are more frequently accessed than others). 
\end{itemize}

\subsection{Symmetric Searchable Encryption}
\label{ch2:sec:symmetric-searchable-encryption}

\subsubsection{Overview}

\gls{SSE} is a type of \gls{SE} based on symmetric key cryptography. We will start by specifying the security definitions used to evaluate such schemes and then describe a few reference schemes.

\subsubsection{Security Definitions}

\gls{IND-CPA} is a property commonly used to evaluate encryption schemes, particularly in the context of public-key or symmetric-key cryptography. Informally, a scheme is \gls{IND-CPA} secure if an adversary $A$, given two arbitrary plaintext messages and their respective ciphertexts, cannot distinguish which ciphertext corresponds to each message with a probability higher than that of random chance. This probability upholds even if $A$ can adaptively query the encryption oracle. 

This definition is not appropriate for \gls{SE} schemes where the leakage of information originates from trapdoors and queries.

In 2006, Curtmola et al. \autocite{Curtmola2011-hi} revisited the security definitions for \gls{SSE} schemes, contributing with two new definitions, which are still the standard in use: \gls{IND-CKA1} and \gls{IND-CKA2}. 

Both require that only the outcome of search patterns may be leaked. Additionally, they guarantee that trapdoors do not leak information about the keywords, apart from what can be inferred from the search and access pattern. Their main difference is whether the adversary can produce an adaptive or non-adaptive attack.

\begin{itemize}
    \item \emph{IND-CKA1}: guarantees security if the client generates all queries at once;
    \item \emph{IND-CKA2}: guarantees security even if the client chooses the queries based on previously obtained trapdoors and search outcomes.
\end{itemize}

\subsubsection{SSE Schemes}

\paragraph{Song et al.} In 2000, the first practical \gls{SSE} scheme was proposed \autocite{Song2000-ir}. The general basis of their method is to encrypt each word of a document and embed hash values with a specific format in the ciphertext. During search, the server executes a \emph{sequential scan} of the document, extracts the hashes and checks for matches on their values. 

First, the plaintext is divided into \emph{fixed-size} blocks $w_i$ and each one is encrypted with a \emph{deterministic}\footnote{Deterministic encryption leaks  message  equality.} encryption algorithm $Enc$, using key $k$. The ciphertexts $X_{i} = Enc_{k}(w_{i})$ are split into two parts $X_{i}=\langle l_{i},r_{i}\rangle$. Using a \emph{stream cipher}\footnote{For information on stream ciphers we recommend reading chapter 8.4 of \autocite{Pearson2017-ve}.} as a \gls{PRF}, a pseudo-random string of bytes $S_{i}=PRF_{k'}(i)$ is generated from the index position using key $k' \neq k$. Then, a hash $Y_{i} = \langle S_{i},F_{k_{i}}(S_{i})\rangle$ of this string of bytes is computed using a keyed hash function $F_{k_{i}}$, where the key is the output of a pseudo-random function $k_{i}=PRF_{k''}(l_{i})$ that takes $l_{i}$ as an input, with key $k''' \neq k''$. Finally, we compute the \emph{XOR} of $X_{i}$ with $S_{i}$ to generate the final ciphertexts $C_{i} = X_{i} \oplus Y_{i}$.

To search, we generate a trapdoor $T = \langle X, K_{s} \rangle$, where $X = E_{k}(w) = \langle L, R \rangle$ is the encryption of keyword $w$ and $K_{s}=PRF_{k''}(L)$. A server searches for the keyword in a document by iterating over all ciphertexts $C_{i}$, computing $C_{i} \oplus X$ and checking if the output is of the form $\langle S, PRF_{K_{s}}(S) \rangle$ for some $S$.

A security analysis of the scheme reveals that it leaks keyword position in the
document. It is also susceptible to statistical attacks (e.g. after various searches, the words in the document can be inferred).


\paragraph{Curtmola et al.} In 2006, two new schemes, for non-adaptive and adaptive adversaries, where proposed \autocite{Curtmola2011-hi}. Each one is, respectively, at least \gls{IND-CKA1} and \gls{IND-CKA2} secure. The idea behind these was to use an \emph{inverted index}, which is an index per distinct word in the database, not per document. 

\emph{Non-Adaptive Scheme} Considering a database with $n$ documents, we define two data structures: an array $A$, where each position of the array contains a node $n_{i,j}$ which belongs to the linked list $L_{i}$ associated to the keyword $w_{i}$; a look-up table $T$ to identify the first node of a linked list associated with a keyword.

To build $A$, we start by constructing the nodes. These nodes are encrypted, stored in random order (using permutations) and augmented with information to find and decrypt the next node in the list. The augmentation step happens before the encryption of the node. Each node $n_{i,j} = \langle id_{m}, k_{next}, p_{next} \, || \, \varnothing \rangle$, includes information about the $id$ of document $m$ ($0 \leq m \leq n$) which contains the keyword, a randomly generated key used to encrypt the next node and a pointer to the next node in the list, or $\varnothing$ if it is the last one. The randomized disposition of the nodes hides the length of the lists.

The look-up table $I$ associates each keyword $w_{i}$ with the node $n_{i,0} = \langle k_{i,1}, p_{i,1} \rangle \oplus PRF_{k}(w_{i})$ stored in a random position using a \gls{PRP}, \linebreak$PRP_{k'}(w_{i})$, where $k \neq k'$. This node contains the encryption key and the pointer to the first node of the linked list $L_{i}$.

To search, we generate trapdoor $T_{w}=\langle PRF_{k}(w), PRP_{k'}(w)\rangle$. The server uses the trapdoor to find the correct node $n_{i,0} = T[PRP_{k'}(w)]$ and then retrieve all the identifiers of the documents that contain the keyword.

\emph{Adaptive Scheme}  Considering a database with $n$ documents, we now define only a look-up table $I$, but instead of generating the index for a single keyword $w$, we generate it for the family of the keyword. The family of a keyword is defined as a set of labels $F_{w} = \{w \, | \, str(j) \, : 1 \leq j \leq t\}$ where $j$ is an integer converted to string, using $str$, $t$ is the number of identifiers of documents containing $w$, and the label is the result of concatenating $w$ with $j$. The look-up table associates each label with the identifier of a document containing $w$, in a random position, using $PRP_{k'}(w_{i} \, | \, str(j))$. To hide differences in the number of labels per keyword, extra entries are padded.

We generate the trapdoor $T_{w}= (t_{1},\;\cdots,\;t{n}) = (PRP_{k}(w \; | \; str(1)),\;\cdots,\;PRP_{k}(w \;|\; str(n)))$ to search. The server uses the trapdoor values to find valid index positions in $I$ and outputs the set of identifiers.

\paragraph{Generalization to the Dynamic Setting} The previous schemes where considered static, as they do not allow modifications on the database after setting up the scheme. In 2012, Kamara et al. \autocite{Kamara2012-lm} proposed a \emph{dynamic} extension to the \emph{non-adaptive} scheme in \autocite{Curtmola2011-hi}. They used a \emph{deletion} array to track the positions on the search array that needed to be modified, after an update. They encrypted the array pointers, using homomorphic encryption, to support modifications on their values without decryption.

\paragraph{Generalization to Structured Data} In 2010, Chase and Kamara \autocite{Chase2010-xx} proposed a new adaptively secure scheme based on the \emph{non-adaptive} scheme in \autocite{Curtmola2011-hi}. This new scheme supports arbitrarily-structured data. \glsxtrlong{STE} is based on generating an inverted index in the form of a padded and permuted dictionary. They associated data items with metadata, allowing information about structure to be preserved if needed. 

\paragraph{Generalization to Graph-Structured Data}
After \gls{STE}, additional research has been conducted to support operations on graph-structured data. 

Some authors extended \gls{STE} to support queries on conceptual graphs \autocite{Poh2012-qs}, but they need to compute all possible query answers \emph{a priori}. 

Others, proposed a \gls{SSE} scheme that supports \emph{subgraph isomorphism queries} \autocite{Cao2011-is}. In their approach, they extract features from the graphs and generate the encrypted-indexes from the extracted features. This scheme was planned for generic graphs.

Finally, in \autocite{Suntaxi2021-sj} they combine an index-based \gls{SSE} with a bucketization algorithm. By doing so, queries do not leak information about the structure of the graph. Although the authors state that their approach works with any kind of query, they only focused on neighbor and adjacency ones. 

\section{Secure Semantic Data Solutions}
\label{ch2:secure-semantic-data}

One of the simplest works in the field focuses solely on protecting individuals in \gls{OWL2} ontologies \autocite{Nejjahi2019-ig}, uniquely encrypting their values and leaving the rest of the data in plaintext. This approach aimed at retaining functionality but leaks too much information. The following solutions try to achieve more robust security guarantees.

\subsection{Self-Enforcing Access Control for Encrypted RDF}

This work researched the implementation of access control in RDF datasets \autocite{Fernandez2017-bh} through the use of functional encryption \autocite{Boneh2010-bw}.

\paragraph{Encryption Scheme} In the proposed encryption scheme, functions correspond to the computation of inner-products over $\mathbb{Z}_N$. The set of possible ciphertext attributes of length $n$ is $\sum = \mathbb{Z}^n_N$ and the class of decryption key functions is $F =\{ f_{\overrightarrow{x}} \; | \; \overrightarrow{x} \in \mathbb{Z}^n_N\}$. An attribute vector $\overrightarrow{y} \in F$ is associated with each ciphertext. A decryption key is a vector $\overrightarrow{x}$ such that $f_{\overrightarrow{x}}(y) = 1 \:\:\: \textrm{\emph{iff}} \:\:\:\sum^n_{i = 1} y_ix_i = 0$

To encrypt an RDF triple $t_i = (s_i,p_i,o_i)$ we functionally encrypt a random seed $m_{t_i}$, with an attribute vector $\overrightarrow{y}_{t_i} = (y_{s_i},y'_{s_i},y_{p_i},y'_{p_i},y_{o_i},y'_{o_i})$ such that $\overrightarrow{y}_l = -r \cdot \sigma(l), \:\: y'_l := r, \:\: l \in \{s,p,o\}$ with $\sigma$ being a function that maps a triple's subject, predicate and object to a value in $\mathbb{Z}_N$, and $r$ a random value also in $\mathbb{Z}_N$

\gls{AES}\footnote{For information on \gls{AES} we recommend reading chapter 6 of \autocite{Pearson2017-ve}.} encryption is used to encrypt the plaintext triple using a key derivable from the generated seed $m_{t_i}$. The final ciphertext $c_{t_i} = \left\langle AES(t_i), \: FE(m_{t_i}) \right\rangle$ corresponds to both the \gls{AES} ciphertext and the ciphertext of the seed.

The decryption keys are generated taking into account the possible triple patterns $tp = (s,p,o)$ combinations. For each triple pattern  we define $\overrightarrow{x}_{tp} =(x_s,x'_s,x_p,x'_p,x_o,x'_o)$ such that $x_l := 1,\: x'_l := \sigma(l) \textrm{ if } l \textrm{ is bound}$ and $x_l := 0,\: x'_l := 0 \textrm{ otherwise}$.

Finally, to decrypt a triple $t = (s_t,p_t,o_t)$, given a triple pattern $tp = (s_{tp},p_{tp},o_{tp})$ we compute the inner product $\overrightarrow{y}_t \cdot \overrightarrow{x}_{tp} = y_{s_t}x_{s_{tp}} + y'_{s_t}x'_{s_{tp}} + y_{p_t}x_{p_{tp}} + y'_{p_t}x'_{p_{tp}} + y_{o_t}x_{o_{tp}} + y'_{o_t}x'_{o_{tp}}$

Only if the result of the inner product is 0 can we decrypt the seed, derive the correct \gls{AES} key and decrypt the triple.

\begin{figure}[h]
    \centering
    \includegraphics[width=0.6\textwidth]{chapter2/3-Index.png}
    \caption{3-Index Indexing}
    \label{fig:3-index}
\end{figure}

\paragraph{Data Structure} The data is stored in a key-value structure where the keys are unique integer identifiers and the values correspond to the encrypted triples. The identifiers are indexed using two strategies.

The first strategy, illustrated in Figure \ref{fig:3-index}, is a \emph{3-Index} approach. We compute secure hashes $h(s)$, $h(p)$ and $h(o)$ for each triple $t = (s, p, o)$ (unbound values are treated as wildcards). Then, we generate three index tables with the permutations of $h(s)$, $h(p)$ and $h(o)$ and associate the permutations with the identifier of the triple.

The second strategy, illustrated in Figure \ref{fig:vertical-partitioning}, is \emph{vertical partitioning}. Instead of indexing the permutations of subjects, predicates and objects. We create one table per predicate, indexing the subject, object pairs that are related via that predicate. Each identifier is indexed by $("so",\;s,\;o)$ and $("os",\;s,\;o)$

\begin{figure}[h]
    \centering
    \includegraphics[width=0.6\textwidth]{chapter2/Vertical-Partitioning.png}
    \caption{Vertical Partitioning Indexing}
    \label{fig:vertical-partitioning}
\end{figure}

\subsection{CryptGraphDB}

\emph{CryptGraphDB} \autocite{Aburawi2018-pq} is an adaptation of the \emph{CryptDB} \autocite{Popa2011-zb, Popa2012-bw, Popa2011-ci} system to graph-structured data. It is capable of answering cypher \autocite{Neo4j_Inc_undated-ft} queries over conceptual graphs. We will start by describing the \emph{CryptDB} architecture, security guarantees, query execution pipeline and encryption schemes. We will then explain the adaptations employed by \emph{CryptGraphDB}.

\paragraph{System Architecture} \emph{CryptDB} is a practical solution for processing \gls{SQL} queries over encrypted relational datasets. This solution leverages onion-encryption to support various primitive operations necessary to answer a \gls{SQL} query (equality checks, order comparisons, aggregates and joins).

The system is composed of an application server, a proxy and an unmodified \gls{DBMS}. The proxy intercepts all queries, encrypts and decrypts all data, and changes query operators while preserving the semantics of the original query. 



\paragraph{Security Threats} \emph{CryptDB} addresses two security threats. The first threat, consists in a passive attacker (e.g. database administrator), which only observes confidential data in the \gls{DBMS} server. Against this attacker \emph{CryptDB} provides a few security guarantees: sensitive data is always encrypted; the \gls{DBMS} server cannot generate results for queries whose primitives do not belong to the subset of supported operations; information leakage depends on the types of computations that will be executed over the data - if no relational predicate filtering is needed, nothing more than the data's length is leaked; on equality checks, it reveals the histogram of encrypted values in the respective columns; in order checks, the system reveals the order between the encrypted values in the column.

The second threat is an attacker that gains control over the proxy, application server and \gls{DBMS} server. Against this adversary, \emph{CryptDB} provides confidentiality of data belonging to users not logged in during the period of the compromise (until it has been fixed). The adversary can, however, delete data stored in the database.

Besides confidentiality, \emph{CryptDB} does not ensure integrity, freshness, or completeness of results.

Figure \ref{fig:cryptdb-system-model} shows the system model for CryptDB and illustrates the scope of the two security threats. In the image, rectangular boxes represent components and processes whereas round boxes represent data. The gray shading represents CryptDB specific elements and the vertical striped lines separate the components, which run in isolated hardware.

\begin{figure}[H]
    \centering
    \includegraphics[width=\textwidth]{chapter2/CryptDB-system-model.png}
    \caption{CryptDB System Model.}
    \label{fig:cryptdb-system-model}
\end{figure}


\paragraph{Query Execution Pipeline} For executing a query issued by the application server, a sequence of steps are necessary: i) the proxy intercepts and rewrites the query, anonymizing the table and column names and encrypting each constant present with the adequate encryption scheme for the desired operation, using a master key; ii) the proxy checks whether it should provide keys to the \gls{DBMS} server, so that adjustments can be done on the encryption layers of the targeted data, if the check is positive it issues an UPDATE query that triggers a custom function to adjust the encryption of the selected columns; iii) the encrypted query is forwarded to the \gls{DBMS} server and is executed using plain \gls{SQL}; iv) the encrypted results are sent to the proxy, which are then decrypted and forwarded to the application server.

\paragraph{Adjustable Encryption Schemes} To avoid unnecessary leakage of information data is encrypted with different layers, like an onion. When operations need encryption schemes with lower security guarantees, the data encryption layers are pealed. For each layer, per onion, per column, the proxy uses the same encryption key, derived from a master key, $MK$. This reduces the leakage of relations between columns and avoids the complexity of having different keys for each row in a column. For table $t$, column $c$, onion $o$, and encryption layer $l$, the proxy would use the encryption key $K_{t,c,o,l} = PRP_{MK}(t,c,o,l)$.

\begin{figure}[H]
    \centering
    \includegraphics[width=\textwidth]{chapter2/CryptDB-onions.png}
    \caption{CryptDB Encryption Onions.}
    \label{fig:cryptdb-onions}
\end{figure}

Figure \ref{fig:cryptdb-onions} provides an overview of the different onions types in \emph{CryptDB} and how different encryption schemes are layered. The supported encryption schemes are:

\begin{itemize}
    \item \emph{Random (RND)}: A probabilistic encryption scheme that provides the most security but does not allow any computations on the ciphertexts.
    \item \emph{Deterministic (DET)}:  A deterministic encryption scheme which allows equality checks to be executed over ciphertexts.
    \item \emph{Order-Preserving (OPE)}: An encryption scheme that preserves order between the ciphertexts in a column. This allows the server to execute range queries.
    \item \emph{Homomorphic encryption (HOM)}: The Paillier cryptosystem \autocite{Paillier1999-js} is used to probabilistically encrypt numeric values. This scheme allows additive operations to be executed with ciphertexts.
    \item \emph{Join (equi-JOIN)}: In \emph{CryptDB} columns are encrypted with different encryption keys. Essentially, joins can not be performed between columns until the encryption keys of the columns coincide. To support join operations between different columns a new primitive is introduced JOIN-ADJ. It can be interpreted as a \emph{keyed random hash} with the property that the hashes can be adjusted to modify their key without access to the plaintext, using an adjustment key $\Delta K = {k_{c_1}/k_{c_2}}$. The adjustment of the columns is done only on one of the two (the first column in lexicographic order). The information leakage is similar to that of the DET layer.
    \item \emph{OPE-JOIN}: \emph{CryptDB} does not provide a construct to dynamically modify the keys of columns encrypted with the \gls{OPE} scheme. To support range joins, these columns have to be specified before, so that they are encrypted under the same key.
    \item \emph{Word search (SEARCH)}: This encryption scheme allows queries with the LIKE keyword. In \emph{CryptDB} they used the scheme by Song et al. \autocite{Song2000-hu}. The information leakage is similar to that of the RND layer.
\end{itemize}

\emph{CryptGraphDB} builds upon the approach of \emph{CryptDB}, adapting the solution to be able to answer \emph{cypher} queries \autocite{Neo4j_Inc_undated-ft}. The proxy is implemented to handle \emph{cypher} queries instead of \gls{SQL} queries, but keeps the same functionality. Furthermore, as the \emph{cypher} language does not require the join operator, \emph{CryptGraphDB} only maintains the RND, DET, OPE, HOM and SEARCH encryption schemes.

\subsection{GOOSE}

\gls{GOOSE} \autocite{Ciucanu2020-lx} is a secure graph outsourcing framework capable of answering \gls{UCRPQ} queries. The solution relies on secure multiparty computation techniques, encryption schemes and a standard SPARQL engine. 

\gls{GOOSE} was designed as a distributed system with 5 participants: a \emph{data owner} (DO), who owns the outsourced graph $G$; a \emph{user} (U), who submits graph queries against $G$ and receives query answers; a \emph{query translator} (QT); \emph{answer translator} (AT); and a \emph{\gls{SPARQL} engine} (SE).

\paragraph{Preliminaries}

Let $\textrm{ENC}_{A\_B}$ or $\textrm{ENC}_{B\_A}$ denote \glsxtrshort{AES-CBC} encryption using a shared key between participants A and B. 

Let the outsourced \gls{RDF} graph be $G = (N, E)$ where $N$ is a set of nodes and $E \in N \times \Sigma \times N$ is a set of directed edges between the nodes, labeled using the alphabet $\Sigma$. Let $\sigma_N: N \to \{0,\:\cdots,\: \vert N \vert - 1\}$ and $\sigma_\Sigma : \Sigma \to \{0,\;\cdots,\; \vert \Sigma \vert - 1\}$ be two random bijections, respectively, for the edges and nodes, and let $\sigma^{-1}$ represent the inverse of $\sigma$. The hidden set of edges from $E$ is defined as $\hat{E} = \{(\sigma_N(s), \sigma_\Sigma(p), \sigma_N(o)) \;\vert\; (s, p, o) \in E\}$.

Consider $\Sigma^+= \{a, a^- \;\vert\; a \in \Sigma \}$, $a^-$ denotes the inverse of the edge label $a$. Let $V = \{?v_0,...,?v_ n\}$, with $n \in \mathbb{Z}^{+}$, be a set of variables. A \emph{query rule} is an expression of the form \emph{head} $\leftarrow$ \emph{body} defined as $(\overline{?v}) \leftarrow (?v_0, r_0, ?v_1),\cdots,(?v_{n+1}, r_n, ?v_{n+2} )$, where, for every $i \in \mathbb{Z}_0^{+}$, $r_i$ is a regular expression over $\Sigma^+$ using $\{\cdot, + , \ast\}$ and $\overline{?v}$ is a vector composed of a subset of the variables from the \emph{body}. The length of this vector is the \emph{arity} of the rule. 

A \gls{UCRPQ} query $q$ is finite and non-empty set of \emph{query rules} with the same \emph{arity}. The answer of a query $q$ evaluated against a graph $G$, using standard \gls{SPARQL} semantics, is denoted as $ans(G, q)$.

\paragraph{Graph Outsourcing} To outsource a graph, the \emph{data owner} generates $\sigma_V$, $\sigma_\Sigma$ and $\hat{E}$. The owner then sends $\textrm{ENC}_{DO\_QT}(\sigma_\Sigma)$ to the \emph{query translator}, $\textrm{ENC}_{DO\_AT}(\sigma_V)$ to the \emph{answer translator} and $\textrm{ENC}_{DO\_SE}(\hat{E})$ to the \gls{SPARQL} engine. 

\paragraph{Query Evaluation} In \gls{GOOSE}, query evaluation is a four-step process: 

\begin{enumerate}
    \item First, the \emph{user} generates a query $q$ and sends $\textrm{ENC}_{U\_QT}(q)$ to the \emph{query translator}. 
    \item The \emph{query translator} generates $\hat{q}$ by applying the $\sigma_\Sigma$ function to all predicates in $q$. $\textrm{ENC}_{QT\_SE}(\hat{q})$ is sent to the \emph{\gls{SPARQL} engine}.
    \item The \emph{\gls{SPARQL} engine} evaluates $\hat{q}$ over $\hat{G} = (\bigcup_{(\hat{s},\hat{p}, \hat{o}) \in \hat{E}} \{\hat{s},\hat{o}\}, \hat{E})$ obtaining the answer $ans(\hat{G}, \hat{q})$. It then sends $\textrm{ENC}_{SE\_AT}(ans(\hat{G}, \hat{q}))$ to the \emph{answer translator}.
    \item The \emph{answer translator} deobfuscates $ans(\hat{G}, \hat{q})$ using ${\sigma_V}^{-1}$ and generates the final answer $ans(G, q)$. It then sends $\textrm{ENC}_{AT\_U}(ans(G, q))$ to the \emph{user}.
\end{enumerate}

Figure \ref{fig:goose} shows both the graph outsourcing (step 0) and the query evaluation (steps 1-4) processes:

\begin{figure}[H]
    \centering
    \includegraphics[width=\textwidth]{chapter2/Goose.png}
    \caption{Graph Outsourcing and Query Evaluation, in \gls{GOOSE}.}
    \label{fig:goose}
\end{figure}

\paragraph{Security Model}
\gls{GOOSE} is designed to be secure against \emph{honest-but-curious} adversaries. It is assumed that data and computations are outsourced to external servers. The solution is expected to have the following security properties, which have been formally proven by the authors: 

\begin{enumerate}
    \item No cloud node can learn $G$.
    \item Let $q$ be a query submitted by a user and $ans(G, q)$ the answer of the query on the graph $G$. No cloud node can learn simultaneously $q$ and $ans(G, q)$.
    \item An external observer cannot learn $G$, any query $q$ nor answer $ans(G, q)$ by analyzing the network messages.
\end{enumerate}

\subsection{Discussion}
\label{ch2:sec:semantic-data:discussion}

\paragraph{Current Solutions Limitations} In this section we covered a few solutions for secure outsourcing of semantic data.

The \emph{CryptGraphDB} \autocite{Aburawi2018-pq} was not designed to support the \emph{SPARQL} query language and suffers from the same security leakage problems inherent to \emph{CryptDB}. Some authors argue that the usage of deterministic encryption to support join operations leaks too much information (e.g. frequency information about the encrypted items) and is vulnerable to \emph{leakage-abuse} attacks \autocite{Jutla2021-af, Hahn2019-sg, Shafieinejad2021-rb}.

The solution from J. D. Fernández et al. \autocite{Fernandez2017-bh} uses hashes as indexes and this leaks the structure of the \gls{RDF} graph. Furthermore, it only supports search operations over simple \emph{triple patterns}.

The \gls{GOOSE} \autocite{Ciucanu2020-lx} solution can answer \gls{SPARQL} queries with complex \emph{graph patterns}, but the data obfuscation methods used are deterministic and still leak the topology of the \gls{RDF} graphs and frequency information of the encrypted keywords.

Furthermore, none of these solutions offer reasoning capabilities, and it is not clear whether they support write operations besides a single initial upload of the data.

\paragraph{Orthogonal Security Research} Like the previous solutions, the main focus of our thesis will be on data \emph{confidentiality}. Nonetheless, we will briefly mention the research that is being done in other security fields regarding \gls{RDF} graphs.  

One of those fields is ensuring \emph{integrity} and \emph{authenticity} of \gls{RDF} data. Various signage schemes \autocite{Carroll2003-zi, Tummarello2005-qb,Kasten2013-rf, Kasten2014-ru, Kasten2014-ny, Sutton2018-ya} and even one service-oriented solution \autocite{Mana2011-bm} have been proposed. 

The other field is \emph{access control}. Most solutions try to embed the policies within the data, using annotations. This allows for high-granularity in the \emph{access control}, which can be inferred, but requires conflict resolution \autocite{Jain2006-zy, Abel2007-sf, Flouris2010-bf, Papakonstantinou2012-cu, Costabello2012-ow,Kirrane2013-ef, Costabello2013-tb}. 


\section{Summary}
\label{ch2:sec:summary}

In Section \ref{ch2:sec:semantic-web}, we described the layered architecture of the \emph{Semantic Web} and discussed the standard languages for describing its \emph{data model} and \emph{ontologies}. More specifically, we covered the semantics and syntax of \gls{RDF} and explained the primitives of \gls{RDFS}. 

Additionally, we specified the structure of \gls{OWL2}, focusing on the mechanism it provides to write ontologies. We also discussed the restriction imposed by \emph{OWL DL} to remain decidable. 

Furthermore, we explained \gls{SPARQL}, providing mathematical definitions, defining its semantics, describing the syntax, enumerating a few examples of queries and reasoning preservation techniques.

We finished Section \ref{ch2:sec:semantic-web} by Listinga few of the tools, software, storage solutions and evaluation benchmarks used in this field of research.

In Section \ref{ch2:sec:secure-computations}, we discussed various solutions to execute secure computations in untrusted servers. We started by describing the \emph{honest-but-curious}, or \emph{passive}, adversary model and explained why \emph{Encryption at Rest} and \gls{ORAM} are not practical approaches. 

We proceeded to cover \glsxtrlong{HE} schemes, informally mentioning the mathematical basis of \gls{HE} and enumerating the type of schemes available, explaining property-preserving encryption schemes and blinding techniques.

In Section \ref{ch2:sec:searchable-encryption}, we explained the difference between \gls{SSE} and \gls{PEKS}. We then described the general model for index-based \gls{SE}, as well as the \emph{Client/Server} architecture of \gls{SE} systems. We continued by discussing how to extend \gls{S/S} architectures, covered the types of privacy leakage.

In Section \ref{ch2:sec:symmetric-searchable-encryption} we specified the standard security definitions in use, and gave an in depth explanation of reference \gls{SSE} schemes. We finished this section by showing recent generalizations, focusing on dynamic solutions and schemes that support structured data.

In Section \ref{ch2:secure-semantic-data}, we mentioned related research on the topic of secure semantic data sharing, discussing the trade-off of each approach. Additionally, we briefly covered parallel work in the fields of data \emph{integrity}, \emph{authenticity} and \emph{access control}. 